The $\RealSeparabilityUp$ dependency lattice also factors through the Bishop regular-real route of \autoref{ch:concrete-instances-bishop-regular-real-namecert}. A separability read that enumerates finite dyadic approximants must first expose the $\RegSeqRatUp$ readback of \autoref{ch:concrete-instances-regseqrat-namecert}, the dyadic tolerance rows of \autoref{ch:concrete-instances-dyadicratcore-namecert}, and the regular-real handoff surface before treating the resulting $\RealUp$ seal as part of a countable dense family.

This keeps the separability route inside displayed L10 data: dyadic density supplies finite approximation addresses, $\RegSeqRatUp$ supplies regular rational names, and $\BishopRegularRealUp$ records the bounded regular real carrier that the separability certificate may consume. No host density theorem, quotient equality of real names, or ambient completion principle is introduced by this child route.
